% !TeX root = ../../Book3.tex
% 纲要标题：### 6.1 Zariski 引理
% 纲要位置：工作记录/计划纲要.md，第 2064 行。
% 纲要细目（写作范围）：
% * 作为域的有限生成代数；域生成与代数生成、有限分母的限制
% * 极大理想的剩余域
% * 与代数闭域假设的联系；闭点扩域后的分裂

\section{Zariski 引理}\label{b3:sec:0601}

把一个有限型代数模掉极大理想，得到的剩余域依然由有限个元素作为代数生成。域的倒数运算与代数生成元的有限性相互约束，结果是这种剩余域只能是基域的有限代数扩张。Noether 正规化使这一事实有一个简洁证明。

\subsection{作为域的有限生成代数}\label{b3:subsec:0601:01}

设 \(k\) 为任意域，\(t\) 在 \(k\) 上超越。\(k(t)\) 虽由一个元素 \(t\) 作为域生成，却必须包含 \(k[t]\) 中每个非零多项式的倒数。有限个有理函数的分母乘积若为 \(D\ne0\)，则它们生成的 \(k\) 代数包含于 \(k[t,1/D]\)，其元素的分母都可写成 \(D\) 的幂。取非恒定多项式 \(1+tD\) 的一个不可约因子 \(h\)，则 \(h\nmid D\)，所以 \(h^{-1}\) 不在这个子环中：否则存在 \(P\in k[t]\) 和整数 \(m\ge0\)，使 \(D^m=hP\)，与 \(h\nmid D\) 矛盾。这说明有限个代数生成元不能提供这个超越方向上所需的全部倒数；下面的引理把这种限制推广到任意有限型域代数。

\begin{lemma}[Zariski]\label{b3:lem:zariski}
设 \(k\) 为域，\(L/k\) 为域扩张，且 \(L\) 作为 \(k\) 代数有限生成，则 \([L:k]<\infty\)。这一结论称为\term[zariski-yinli]{Zariski 引理}[Zariski's lemma]。
\end{lemma}
\begin{proof}
由\cref{b3:thm:noether-normalization}，存在多项式子环 \(R=k[y_1,\ldots,y_d]\subseteq L\)，使 \(L\) 为有限 \(R\) 代数，因而在 \(R\) 上整。因为 \(L\) 是域，\cref{b3:lem:integral-field-criterion}迫使 \(R\) 也是域。若 \(d>0\)，变量 \(y_1\) 在多项式环中非零且不可逆，矛盾。因此 \(d=0\)，\(R=k\)，而正规化结论正说明 \(L\) 是有限生成的 \(k\) 模，即有限维 \(k\) 向量空间。
\end{proof}

反过来，有限代数扩张选一个有限的向量空间基即可作为代数生成集，所以引理刻画了哪些域是有限型 \(k\) 代数。

\subsection{极大理想的剩余域}\label{b3:subsec:0601:02}

\begin{corollary}\label{b3:cor:finite-type-residue-fields}
设 \(k\) 为域，\(A\) 为交换含幺有限型 \(k\) 代数。对 \(A\) 的任一极大理想 \(\mathfrak m\)，剩余域 \(A/\mathfrak m\) 是 \(k\) 的有限扩张。对 \(A\) 的一般素理想 \(\mathfrak p\)，剩余域 \(\kappa(\mathfrak p)=\operatorname{Frac}(A/\mathfrak p)\) 则只必然是有限生成的域扩张。
\end{corollary}
\begin{proof}
\(A/\mathfrak m\) 是域，同时是 \(A\) 的有限型代数商，应用\cref{b3:lem:zariski}即可。对一般素理想，\(A/\mathfrak p\) 的有限组代数生成元也作为域生成元生成其分式域，但此时已经允许把任意非零多项式倒置，不能继续应用 Zariski 引理。
\end{proof}

例如 \(k[x,y]\) 的极大理想若取 \((x,y)\)，剩余域为 \(k\)；素理想 \((x)\) 的剩余域却是 \(k(y)\)。闭点和一般点的差别，在这里反映为剩余域的代数性与超越性。

\subsection{代数闭域假设的作用}\label{b3:subsec:0601:03}

若 \(k\) 代数闭，它没有非平凡有限代数扩张：有限扩张中的每个元素有 \(k\) 上的不可约最小多项式，而代数闭性迫使该多项式为一次。因此\cref{b3:cor:finite-type-residue-fields}给出
\[
A/\mathfrak m=k
\]
的典范 \(k\) 代数同构。这里的等号意为结构映射 \(k\to A/\mathfrak m\) 为同构，而不是选定某个任意的抽象域同构。

\ProseParagraphBreak
基域不代数闭时，有限剩余域扩张必须保留。例如 \(\mathbb R[x]/(x^2+1)\cong\mathbb C\)，对应一个实多项式环的极大理想，却不对应实数轴上的点。把系数域扩大到 \(\mathbb C\) 后，\(x^2+1=(x-i)(x+i)\)，两因子互素，中国剩余定理给出
\[
\mathbb C[x]/(x^2+1)\cong\mathbb C\times\mathbb C.
\]
原来的一个闭点由此分裂成 \(i,-i\) 两个复坐标点。另一方面，\(\mathbb F_q[x]\) 的次数 \(r\) 的首一不可约多项式给出剩余域 \(\mathbb F_{q^r}\)。下一节正是把代数闭性用在排除这些额外的剩余域上。
